\documentclass[../main.tex]{subfiles}
\begin{document}

\chapter{メニーコア}
\lessoninfo{
  video={Lesson 22，動画 1--19},
  textbook={H\&P \cite{hennessy2017} §1.5，§5.4，付録 F},
  keywords={メニーコアプロセッサ，コヒーレンストラフィック，ネットワークオンチップ，メッシュ，トーラス，タイル，オフチップトラフィック，分散最終レベルキャッシュ，オンチップディレクトリ，部分ディレクトリ，電力バジェット，ターボモード，論理プロセッサ},
}

第18章から第21章では，コアがふつうバスと最終レベルキャッシュを共有し，
バス上のスヌーピングで（共通のバスを持たないマシンではディレクトリで）
キャッシュをコヒーレントに保ち，ロックで協調し，共有データについては一貫性
モデルに頼るマルチプロセッサを構築してきた．しかし世代が進むたびに1つの
チップにはより多くのコアが収まるようになり，4個や8個のコアをうまく扱える
機構も，数十個から数百個のコアを持つチップでは破綻する．この章では，
そのようなチップが直面する課題---コヒーレンストラフィック，オフチップ
トラフィック，コヒーレンスディレクトリの大きさ，共有された電力バジェット，
そしてオペレーティングシステムによるスレッドの配置---と，講義がそれぞれに
与える解決策を調べる．

\section{コヒーレンストラフィック}
\label{sec:l22-coherence}

\term[めにーこあぷろせっさ]{メニーコアプロセッサ}[many-core processor]%
とは，1つのチップに多数のコアを載せたマルチコアプロセッサ，すなわち2個や
8個ではなく数十個から数百個のコアを持つものである．\source{Lesson 22，
動画 1--2；H\&P 第 5 章．}このようなチップの第1の課題はキャッシュ
コヒーレンス（第19章）に関わる．他のコアがキャッシュしている共有データへの
書き込みはすべて\emph{コヒーレンストラフィック}を生む．すなわち，他の
キャッシュへ送られる無効化と，それらのコアが再びそのブロックにアクセス
するときのコヒーレンスミスである．コアが増えればチップ全体では毎秒の
書き込みが増えるので，コヒーレンスメッセージの発生率もコア数とともに増える．
共有バスでは，このトラフィックがまもなくボトルネックになる．
\begin{itemize}
  \item バスは一度に1つの転送しか運べないので，コアは互いのコヒーレンス
    メッセージとキャッシュミスを待つことになる．
  \item より多くのコアをつなぐバスは長くなり，接続されるデバイスも増える
    ので，1回の転送自体も遅くなる．
\end{itemize}

したがってメニーコアチップには2つのものが必要である．コア数とともに容量が
増えるオンチップの相互接続網と，1本のバス上のすべての操作をすべての
キャッシュが観測することに依存しないコヒーレンス機構である．後者が
ディレクトリ方式（\cref{sec:l19-directory}）であり，これは各メッセージを
すべてのコアへブロードキャストするのではなく，そのブロックを保持する
キャッシュだけへ届ける．

\section{ネットワークオンチップ}
\label{sec:l22-noc}

共有バスはネットワークに置き換えられる．\source{Lesson 22，動画 3--4；
H\&P 付録 F；\cite{dally2001}．}

\begin{definition}[ネットワークオンチップ]
  \term[ねっとわーくおんちっぷ]{ネットワークオンチップ}[network on chip]%
  は，チップのコアを，コアごとに1つ置かれたルータの間の1対1のリンクで
  接続する．メッセージは宛先に着くまでルータからルータへと渡り，異なる
  リンク上のメッセージは同時に進む．
\end{definition}

\term[めっしゅ]{メッシュ}[mesh]トポロジでは，ルータが2次元の格子をなし，
各ルータは東西南北の隣接ルータとリンクで結ばれる
（\cref{fig:l22-topologies}b）．どのコアも，いくつかのリンク，すなわち
\emph{ホップ}からなる経路を通って他のどのコアにも到達し，異なるリンクを
使う転送は同時に進むので，ネットワーク全体では1本のリンクの何倍もの
トラフィックを運ぶ．さらに2つの性質がメッシュをチップに適したものにする．
\begin{description}
  \item[スケーリング] コアを増やすと格子が大きくなり，追加されたコアが
    リンクも追加するので，ネットワークの容量はバスと違ってコア数とともに
    増える．どのリンクも隣り合う2つのコアを結ぶだけなので，チップがどれ
    ほど大きくなっても短いままであるのに対し，バスはすべてのコアに
    またがらなければならない．
  \item[平面的な配置] リンクは隣接するもの同士だけを結び，交差しないので，
    メッシュはチップの平らな表面に容易に配置できる．
\end{description}
このようなチップは\term[たいる]{タイル}[tile]を複製して作られる．タイル
とは，コアとそのプライベートキャッシュ，そしてそのルータであり，以降の節
では全コアが共有する構造の一部分もここに加わる．

\begin{figure}[htbp]
  \centering
  % Interconnect topologies for many cores: (a) a shared bus, (b) a 4x4 mesh
% carrying two transfers at the same time on different links, (c) a 4x4
% torus, i.e. a mesh with wraparound links in every row and column.
% Links are drawn between node centres on the background layer, so the
% filled nodes hide the parts of the lines inside them.
\begin{tikzpicture}[x=1cm,y=1cm, font=\footnotesize\sffamily,
  nd/.style={draw, fill=ln-accent!15, minimum size=0.36cm, inner sep=0pt},
  lk/.style={ln-muted, line width=0.5pt},
  wrap/.style={ln-accent, line width=0.5pt, rounded corners=2pt},
  route/.style={line width=1.6pt, line cap=round, line join=round},
  paneltitle/.style={font=\footnotesize\sffamily, anchor=north}]
  % ---------------------------------------------------------------- (a) bus
  \begin{scope}
    \foreach \i in {0,...,3} {
      \node[nd] at (0.3+\i*0.8,2.3) {};
      \node[nd] at (0.3+\i*0.8,0.4) {};
    }
    \begin{scope}[on background layer]
      \foreach \i in {0,...,3}
        \draw[lk] (0.3+\i*0.8,2.3) -- (0.3+\i*0.8,0.4);
    \end{scope}
    \draw[line width=1.4pt] (-0.1,1.35) -- (3.1,1.35);
    \node[paneltitle] at (1.5,-0.6) {\iflangja{(a) バス}{(a) bus}};
  \end{scope}
  % --------------------------------------------------------------- (b) mesh
  \begin{scope}[xshift=4.0cm]
    \foreach \x in {0,...,3} \foreach \y in {0,...,3}
      \node[nd] at (\x*0.9,\y*0.9) {};
    \begin{scope}[on background layer]
      \foreach \k in {0,...,3} {
        \draw[lk] (\k*0.9,0) -- (\k*0.9,2.7);
        \draw[lk] (0,\k*0.9) -- (2.7,\k*0.9);
      }
      % two transfers in progress on disjoint links
      \draw[route, ln-caution] (0,2.7) -- (1.8,2.7) -- (1.8,1.8);
      \draw[route, ln-tip] (2.7,0) -- (2.7,0.9) -- (0.9,0.9);
    \end{scope}
    \node[paneltitle] at (1.35,-0.6) {\iflangja{(b) メッシュ}{(b) mesh}};
  \end{scope}
  % -------------------------------------------------------------- (c) torus
  \begin{scope}[xshift=8.1cm]
    \foreach \x in {0,...,3} \foreach \y in {0,...,3}
      \node[nd] at (\x*0.9,\y*0.9) {};
    \begin{scope}[on background layer]
      \foreach \k in {0,...,3} {
        \draw[lk] (\k*0.9,0) -- (\k*0.9,2.7);
        \draw[lk] (0,\k*0.9) -- (2.7,\k*0.9);
        % wraparound links: below each row, right of each column
        \draw[wrap] (0,\k*0.9) -- (-0.32,\k*0.9) -- (-0.32,\k*0.9-0.3)
          -- (3.02,\k*0.9-0.3) -- (3.02,\k*0.9) -- (2.7,\k*0.9);
        \draw[wrap] (\k*0.9,2.7) -- (\k*0.9,3.02) -- (\k*0.9+0.3,3.02)
          -- (\k*0.9+0.3,-0.32) -- (\k*0.9,-0.32) -- (\k*0.9,0);
      }
    \end{scope}
    \node[paneltitle] at (1.35,-0.6) {\iflangja{(c) トーラス}{(c) torus}};
  \end{scope}
\end{tikzpicture}

  \caption{メニーコアのための相互接続網．バスは一度に1つの転送しか運ばない．
    メッシュは強調した2つの転送を異なるリンク上で同時に運ぶ．トーラスは
    各行各列の最初と最後のルータの間に折り返しリンクを加える---図では
    左端の列のものと最下行のものだけを描いている．}
  \label{fig:l22-topologies}
\end{figure}

経路を短くする他のトポロジもあるが，配置が難しくなるという代償を払う．
\term[とーらす]{トーラス}[torus]は，各行各列の最初と最後のルータの間に
折り返しリンクを加えたメッシュであり（\cref{fig:l22-topologies}c），最長
経路をほぼ半分にする．ただし折り返しリンクは長いか，他のリンクと交差する．
\emph{フラッテンドバタフライ}と\emph{ハイパーキューブ}は，各ルータをより
多くのルータに結ぶ---フラッテンドバタフライでは自分と同じ行および列の
すべてのルータへ，ハイパーキューブでは $\log_{2}N$ 個のルータへ---ので
経路はさらに短くなるが，長く交差するリンクを必要とし，チップ上に作るのは
難しい．

トラフィックがリンクの上に均等に広がるとすれば，それぞれ毎秒 $C$ 個まで
メッセージを運ぶ $L$ 本のリンクからなるネットワークは，毎秒
\begin{equation}
  \text{スループット} \;\approx\; \frac{L\cdot C}{h}
  \label{eq:l22-throughput}
\end{equation}
個のメッセージを運ぶ．ここで $h$ はメッセージ当たりの平均ホップ数であり，
メッセージはホップごとに1本のリンクを占有するからである．$k\times k$ の
メッシュはリンクを $L=2k(k-1)$ 本持ち，宛先を他のタイルから一様に選ぶ場合
の平均距離は $h=2k/3$ ホップで，これはコア数 $N=k^{2}$ の平方根でしか
増えない．

\cref{eq:l22-throughput} は平均負荷による見積もりであり，現実のルーティング
がこれに達することはない．1本のリンクの負荷はそこを経由するすべての
メッセージの発生率の和であり，ネットワークが与えられたトラフィックを運べる
のはどのリンクも $C$ を下回っている間だけなので，限界を決めるのは平均では
なく最も混雑したリンクである．\caution{メッシュはスループットを上げるが，
レイテンシを短くはしない．遠いタイルへのメッセージは多くのルータを通り，
そのそれぞれが遅延を加える．}メッセージの宛先がチップ全体に散らばり，
ルータが各メッセージをまず一方の次元に沿って，次にもう一方の次元に沿って
送る場合，メッシュの中央を横切るリンクが最も混雑し，それらは毎秒およそ
$2kC$ 個のメッセージで飽和する．したがってこのようなトラフィックに対する
メッシュの容量は $\sqrt{N}$ で増え，$N$ に比例して増えるのはメッセージが
隣り合うタイルの間に留まる場合だけである．対照的にバスは，いくつのコアを
つなごうともリンク1本分を超えるトラフィックを運ぶことはない．トラフィック
がチップを横切るようになるとコア当たりの容量は $1/\sqrt{N}$ で下がるので，
コアのデータをそのコア自身のタイルに置くことが重要になる
（\cref{sec:l22-llc}）．

\begin{example}
  \label{ex:l22-throughput}
  各リンクとバスが毎秒 $C$ 個のメッセージを運ぶとする．16 コアでは，バスは
  $C$ を運ぶ．$4\times4$ のメッシュはリンクを 24 本持ち，平均距離は
  $8/3\approx 2.67$ ホップなので，\cref{eq:l22-throughput} はこれを
  $24C/2.67\approx 9C$ と見積もるが，中央のリンクがおよそ $8C$ に制限する．
  64 コアでも，バスが運ぶのは高々 $C$ である（バスが長くなる分，実際には
  それ以下になる）．一方 $8\times8$ のメッシュはリンクを 112 本持ち，平均
  距離は $16/3\approx 5.33$ ホップで，見積もりは $21C$，中央のリンクによる
  限界はおよそ $16C$ となる．コア数は4倍だが容量は2倍である．すべての
  トラフィックが隣り合うタイルへ向かうなら（$h=1$），2つのメッシュは
  $24C$ と $112C$ を運ぶことになる．
\end{example}

\section{オフチップトラフィックと分散 LLC}
\label{sec:l22-llc}

第2の課題はチップの外へ出るトラフィックである．\source{Lesson 22，動画
5--7．}各コアは自分のキャッシュを持ち，コア当たりのミス率はコア数によって
変わらないので，コアが多いチップほど主記憶へ比例して多くの要求を送る．
これらの要求はチップのピンを通る．ピンの数も世代ごとに増えるが，コア数より
はるかにゆっくりとしか増えないので，メモリインタフェースがボトルネックに
なる．したがって\emph{コア当たり}のメモリ要求数を減らさなければならない．
より多くのアクセスをチップ上で，すなわち全コアが共有し，その容量がコア数に
おおよそ比例して増える最終レベルキャッシュで処理しなければならない．しかし
1つの大きな LLC は遅い．キャッシュが大きいほどヒット時間は長くなり
（第12章），さらに全コアの L2 ミスが通らなければならない入口が1か所しか
なく，それ自体がボトルネックになるからである．

\begin{definition}[分散最終レベルキャッシュ]
  \term[ぶんさんさいしゅうれべるきゃっしゅ]{分散最終レベルキャッシュ}%
  [distributed last-level cache]は，論理的には全コアが共有する1つの
  キャッシュ---各ブロックはコアごとに1回ではなく，高々1回だけ格納される
  ---でありながら，物理的には\emph{スライス}に分割され，各タイルに1つの
  スライスが置かれる（\cref{fig:l22-llc}）．
\end{definition}

LLC の容量はコア数とスライス1つの容量の積なので，コア数とともに自動的に
増える一方，各スライスは小さく高速なままであり，コアの要求は多数の
スライスに分散される．

\begin{figure}[htbp]
  \centering
  % A many-core chip of eight tiles connected by a mesh.  Each tile holds a
% core, its private L1 and L2 caches, one slice of the distributed LLC, the
% part of the on-chip directory kept in that tile, and a router (R).  An L2
% miss of core 1 on a block of a set whose number ends in 110 goes to LLC
% slice 6.
\begin{tikzpicture}[x=1cm,y=1cm, font=\scriptsize\sffamily, >={Stealth[length=1.8mm]},
  tile/.style={draw=ln-rule, rounded corners=2pt},
  core/.style={draw, fill=ln-accent!15, minimum width=1.5cm, minimum height=0.34cm, inner sep=0pt},
  box/.style={draw, fill=ln-box, minimum height=0.34cm, inner sep=0pt},
  rt/.style={draw, circle, fill=white, minimum size=0.36cm, inner sep=0pt, font=\tiny\sffamily},
  lk/.style={ln-muted, line width=0.6pt},
  hot/.style={ln-caution, line width=1.5pt}]
  \foreach \i in {0,...,7} {
    \pgfmathtruncatemacro{\col}{mod(\i,4)}
    \pgfmathtruncatemacro{\row}{div(\i,4)}
    \begin{scope}[shift={(\col*2.65,-\row*2.6)}]
      \draw[tile] (-1.1,-1.35) rectangle (1.2,0.95);
      \node[core] (core\i) at (-0.2,0.62) {\iflangja{コア}{core}~\i};
      \node[box, minimum width=0.72cm] (l1\i) at (-0.59,0.18) {L1};
      \node[box, minimum width=0.72cm] (l2\i) at (0.19,0.18) {L2};
      \node[box, minimum width=1.5cm] (llc\i) at (-0.2,-0.26) {LLC~\i};
      \node[box, minimum width=1.5cm] (dir\i) at (-0.2,-0.70) {\iflangja{ディレクトリ}{directory}~\i};
      \coordinate (rp\i) at (0.95,-1.12);
    \end{scope}
  }
  % mesh links between routers (drawn before the routers)
  \foreach \i/\j in {0/1,1/2,2/3,4/5,5/6,6/7,0/4,1/5,2/6,3/7}
    \draw[lk] (rp\i) -- (rp\j);
  % route of the request: R1 -> R2 -> R6
  \draw[hot] (rp1) -- (rp2) -- (rp6);
  \fill[ln-caution!20] (llc6.south west) rectangle (llc6.north east);
  \draw (llc6.south west) rectangle (llc6.north east);
  \node at (llc6) {LLC~6};
  \foreach \i in {0,...,7} \node[rt] (r\i) at (rp\i) {R};
  \draw[->, hot, line width=0.9pt] (l21.east) -| (r1.north);
  \draw[->, hot, line width=0.9pt] (r6.west) -| ([xshift=0.2cm]llc6.east) -- (llc6.east);
\end{tikzpicture}

  \caption{メッシュで接続された8つのタイル．各タイルはコア，その L1 と L2
    キャッシュ，分散 LLC の1スライス，オンチップディレクトリの一部
    （\cref{sec:l22-directory}），そしてルータ（R）を保持する．セット番号が
    2進で 110 で終わるブロックに対するコア 1 の L2 ミスは，LLC スライス 6
    へ送られる．}
  \label{fig:l22-llc}
\end{figure}

L2 ミスが起きると，コアは要求を，そのブロックを保持しうる唯一のスライスへ
送らなければならず，そのスライスをアドレスだけから計算できなければならない．
講義はブロックをタイルにラウンドロビンで分配する方法を2通り挙げる．
\begin{description}
  \item[キャッシュインデックスによる] セット番号の下位ビットがタイルを
    選ぶので，連続するセットは連続するタイルへ行く．ブロックは均等に
    分散されるが，どのコアがそれを使うかは考慮されない．あるスレッドの
    ブロックは全タイルに散らばり，その LLC アクセスの大半は遠いタイルへ
    向かう．
  \item[ページ番号による] 物理ページ番号のビットがタイルを選ぶので，
    1つのページのブロックはすべて同じスライスに入る．各ページがどの物理
    フレームを受け取るかはオペレーティングシステムが決めるので，
    スレッドには，そのスライスがスレッド自身のタイルかその近くにある
    フレームを与えることができる．こうすれば LLC アクセスの大半はローカル
    に留まる---NUMA マシンでスレッドのメモリをそのノードに置くのと同じ
    考え方である（\cref{sec:l18-distributed}）．
\end{description}

\begin{example}
  \label{ex:l22-llc}
  あるチップは \SI{1}{MB} の LLC スライスを持つタイルを 32 個備え，その
  LLC は \SI{32}{MB} を保持する．\SI{64}{B} のブロックと 8 ウェイセット
  アソシアティブ構成なら，LLC のセット数は
  $2^{25}/(2^{6}\cdot 2^{3})=2^{16}$ であり，各スライスはそのうち
  $2^{11}=2048$ を保持する．タイルをキャッシュインデックスで選ぶ場合，
  16 ビットのセット番号はアドレスのビット 6--21 からなり，その下位 5
  ビット（アドレスのビット 6--10）が 32 個のタイルの1つを選び，残りの
  11 ビット（アドレスのビット 11--21）がそのスライス内のセットを選ぶ．
  タグは，$2^{16}$ セットを持つどのキャッシュとも同じく，ビット 22 以上の
  アドレスビットからなる．
\end{example}

\section{オンチップディレクトリ}
\label{sec:l22-directory}

ディレクトリ方式にも固有の課題がある．\source{Lesson 22，動画 8--11；
H\&P §5.4；\cite{agarwal1988}．}\cref{sec:l19-directory} のディレクトリは
メモリブロックごとにエントリを持ち，各エントリはコアごとの存在ビットを
1つずつと排他ビットを1つ持つ．数十ギガバイトの主記憶は数億のブロックから
なり，そのすべてにエントリを持つディレクトリはチップに載らない．

\begin{example}
  \label{ex:l22-directory-size}
  64 コアのシステムが \SI{32}{GB} のメモリを \SI{64}{B} のブロックに
  分けて持つとすると，ブロック数は $2^{35}/2^{6}=2^{29}\approx 5.37$ 億で
  ある．1エントリ 65 ビットとすると，完全なディレクトリにはおよそ
  $3.5\times10^{10}$ ビット，すなわち \SI{4.1}{GB} が必要であり，これは
  主記憶自体の約 \SI{13}{\percent} に当たる．
\end{example}

2つの決定がディレクトリをチップに適合させる．第1はエントリをどこに置くか
である．ディレクトリはタイルに分割され，あるブロックのエントリは，その
ブロックの LLC スライスを保持するタイル，すなわちそのブロックの
\emph{ホーム}タイルに置かれる．プライベートキャッシュがミスしたコアは，
LLC を探すためにいずれにせよそのタイルに問い合わせるので，追加の
メッセージなしにディレクトリエントリに到達する．また，どのタイルも中央の
ボトルネックにならない．第2の決定，すなわちディレクトリを十分小さくする
決定は，そもそもどのブロックがエントリを持つかである．

\begin{definition}[オンチップディレクトリ]
  \term[おんちっぷでぃれくとり]{オンチップディレクトリ}[on-chip directory]%
  とは，チップ上に置かれるコヒーレンスディレクトリである．これはタイルに
  分散され，各ブロックのエントリはそのブロックの LLC スライスのタイルに
  置かれる．さらにこれは，タイルごとに固定個数のエントリを持ち，一部の
  ブロックについてだけエントリを保持する
  \term[ぶぶんでぃれくとり]{部分ディレクトリ}[partial directory]である．
\end{definition}

エントリが必要なのは，少なくとも1つのプライベートキャッシュ（L1 または
L2）が保持しているブロック，すなわちエントリの存在ビットが少なくとも1つ
立つブロックだけである．LLC か主記憶にしかないブロックのエントリはすべて
0 になるが，エントリが無いことも同じことを述べている．そこでディレクトリ
は，エントリを持たないブロックをプライベートキャッシュが取得したときに
エントリを割り当てる．そして，どのプライベートキャッシュもそのブロックを
保持しなくなれば，そのエントリは不要になる．新しいエントリが必要なのに
そのタイルのエントリがすべて使われているとき，ディレクトリは1つを
置き換える．
\begin{enumerate}
  \item 犠牲とするエントリ，例えば最も長く使われていないものを選ぶ．
  \item 犠牲エントリで存在ビットが立っているすべてのタイルへ無効化を送る．
    これらのキャッシュはそのブロックのコピーを破棄する（変更されたコピーは
    先に書き戻される）．
  \item そのエントリを新しいブロックに使う．
\end{enumerate}
無効化されたキャッシュは，他のどのコアもそのブロックに書いていないにも
かかわらず，次にそのブロックへアクセスするときにミスする．このように
ディレクトリの置換は，3C モデル（第14章）のミスと第19章のコヒーレンスミス
に加えて，新しい種類のキャッシュミスを生む．ディレクトリのエントリが多い
ほど，このようなミスは稀になる．\margin{部分ディレクトリは，独自の置換
アルゴリズムを持つディレクトリエントリのキャッシュのように振る舞う．}

\begin{example}
  \label{ex:l22-directory-trace}
  \cref{tab:l22-directory-trace} は，C0--C3 の4つのコアを持つチップで，
  エントリを2つ置ける1つのタイルのディレクトリ部分を追跡したものである．
  ブロック A，B，C はいずれもこのタイルをホームタイルとし，初めはどの
  プライベートキャッシュもこれらを保持していない．
\end{example}

\begin{table}[htbp]
  \centering\small
  \caption{エントリを2つ持つ部分ディレクトリ．各エントリはブロック，
    排他ビット，C0--C3 の存在ビットを示す．}
  \label{tab:l22-directory-trace}
  \begin{tabular}{@{}clccl@{}}
    \toprule
    手順 & 事象 & エントリ 1 & エントリ 2 & ディレクトリの動作 \\
    \midrule
    1 & C0 が A を読む  & \texttt{A 1 1000} & ---               & エントリ 1 を割り当て \\
    2 & C2 が A を読む  & \texttt{A 0 1010} & ---               & 読み出しを C0 へ転送 \\
    3 & C1 が B に書く  & \texttt{A 0 1010} & \texttt{B 1 0100} & エントリ 2 を割り当て \\
    4 & C3 が C を読む  & \texttt{C 1 0001} & \texttt{B 1 0100} & A を置換: C0，C2 を無効化 \\
    5 & C0 が A を読む  & \texttt{C 1 0001} & \texttt{A 1 1000} & B を置換: C1 を無効化 \\
    \bottomrule
  \end{tabular}
\end{table}

手順 4 ではディレクトリが満杯であり，手順 2 で最後に使われた A が最も長く
使われていないエントリなので，C0 と C2 は A のコピーを失う．手順 5 の C0
による A の読み出しは，この置換だけが原因のミスであり，それがさらに B を
ディレクトリから，そして C1 から追い出す．C1 の変更されたコピーは LLC へ
書き戻される．

\begin{keypoint}
  分散 LLC もオンチップディレクトリも，共有される構造をタイルに分割する
  ことで，中央のボトルネックを作らずにコア数とともに成長させる．ブロックの
  LLC スライスとそのディレクトリエントリは同じタイルに置かれ，ディレクトリ
  はプライベートキャッシュにあるブロックについてだけエントリを保持する
  ---その代償として，ディレクトリの置換によるミスがときおり生じる．
\end{keypoint}

\section{電力バジェットの共有}
\label{sec:l22-power}

第3の課題は電力である．\source{Lesson 22，動画 12--14；H\&P §1.5．}チップ
が消費できる電力は，そのパッケージと冷却が取り除ける分だけであり，この電力
バジェットはすべてのコアで共有される．コアが多いチップほど各コアが受け取る
電力は小さくなり，その電源電圧とクロック周波数も低くしなければならない．
コアの動的電力は $V^{2}f$ に比例し（\cref{eq:active-power}），コアが到達
できる周波数はおおよそ電圧に比例するので，電力はおおよそ $f^{3}$ で増える
（\cref{ex:l18-power}）．$N$ 個のコアが1個分のバジェットを分け合うなら，
各コアはその $1/N$ を受け取り，
\begin{equation}
  f_{N} \;=\; \frac{f_{1}}{\sqrt[3]{N}}
  \label{eq:l22-freq}
\end{equation}
で動作する．ここで $f_{1}$ は単一コアの周波数である．この関係は，電圧を
周波数に比例して下げられることを前提としている．電圧がその下限に近づくと
電力は $f$ に対して線形にしか下がらず，バジェットの分割は周波数の面で
さらに大きな代償を伴う．

\begin{marginfigure}
  \resizebox{\linewidth}{!}{% Performance of N cores sharing the power budget of one core, relative to
% the single core (P proportional to f^3, so each core runs at N^(-1/3) of
% the single-core frequency).  Curves: a program that is fully parallel,
% one with parallel fraction 0.9, and a single-threaded program.
\begin{tikzpicture}[x=0.75cm,y=0.8cm, font=\scriptsize\sffamily]
  \draw[ln-rule] (0,0) grid[xstep=1,ystep=1] (6,4);
  \draw[->] (0,0) -- (6.3,0);
  \draw[->] (0,0) -- (0,4.3);
  \foreach \t/\n in {0/1,1/2,2/4,3/8,4/16,5/32,6/64}
    \node[below, inner sep=1.5pt] at (\t,0) {\n};
  \foreach \v in {0,...,4}
    \node[left, inner sep=1.5pt] at (0,\v) {\v};
  \node[below] at (3,-0.45) {\iflangja{コア数 $N$}{cores $N$}};
  \node[rotate=90, above] at (-0.45,2) {\iflangja{相対性能}{relative performance}};
  \draw[thick, ln-muted, dashed] plot[domain=0:3, samples=30, smooth]
    (\x, {2^(2*\x/3)});
  \node[ln-muted, right, inner sep=1pt] at (3.15,3.7) {\iflangja{全並列}{fully parallel}};
  \draw[thick, ln-accent] plot[domain=0:6, samples=60, smooth]
    (\x, {2^(-\x/3)/(0.1+0.9/2^\x)});
  \node[ln-accent, anchor=north, inner sep=1pt] at (5.0,1.95) {$f_{\mathrm{par}}=0.9$};
  \draw[thick, ln-caution] plot[domain=0:6, samples=40, smooth]
    (\x, {2^(-\x/3)});
  \node[ln-caution, anchor=south west, inner sep=1pt] at (0.1,0.08) {\iflangja{単一スレッド}{single-threaded}};
\end{tikzpicture}
}
  \caption{1コア分の電力バジェットを $N$ 個のコアが分け合うときの性能．
    \cref{eq:l22-perf} による．}
  \label{fig:l22-power-perf}
\end{marginfigure}
CPI は変わらないので，単一スレッドのプログラムは，同じバジェットの単一
コアで走らせる場合に比べて $N$ 個のコアを持つチップでは $\sqrt[3]{N}$ 倍
遅くなる．コアを増やすと逐次プログラムは遅くなるのである．並列部分の割合が
$f_{\mathrm{par}}$ であるプログラムは，Amdahl の法則
（\cref{eq:l18-amdahl-cores}）が述べるように増えたコアから利益を得るが，
周波数の低下によって損をする．単一コアを基準とすると，
\begin{equation}
  \text{性能}(N) \;=\; \frac{1}{\sqrt[3]{N}}\cdot
  \frac{1}{(1-f_{\mathrm{par}}) + \dfrac{f_{\mathrm{par}}}{N}}
  \label{eq:l22-perf}
\end{equation}
である．より一般に，単一コアでの実行時間のうち割合 $t_{p}$ が $p$ 個の
スレッドを使えるなら，そのプログラムは $N$ 個のコアではその時間の
$\sqrt[3]{N}\,\sum_{p} t_{p}/\min(p,N)$ を要する．あるコア数を超えると，
周波数の損失が並列性の利得を上回る（\cref{fig:l22-power-perf}）．

\begin{example}
  \label{ex:l22-power}
  あるチップのバジェットは \SI{120}{W} であり，それをすべて使う単一コアは
  \SI{4.2}{GHz} で動作する．8個のコアはそれぞれ \SI{15}{W} を受け取り，
  $\SI{4.2}{GHz}/\sqrt[3]{8}=\SI{2.1}{GHz}$ で動作するので，単一スレッド
  のプログラムは単一コアでの半分の速度で走る．$f_{\mathrm{par}}=0.9$ の
  プログラムは単一コアの $0.5/(0.1+0.9/8)\approx 2.35$ 倍の性能を示す．
  27 個のコアを \SI{1.4}{GHz} で動かせば $(1/3)/(0.1+0.9/27)=2.5$ に
  達するが，64 個を \SI{1.05}{GHz} で動かすと
  $0.25/(0.1+0.9/64)\approx 2.19$ にしかならない．
\end{example}

\section{ターボモード}
\label{sec:l22-turbo}

プログラムに並列性が無いときは，遊んでいるコアの電力バジェットを，その
プログラムを実行するコアに与えることができる．\source{Lesson 22，動画 15；
H\&P §1.5．}\term[たーぼもーど]{ターボモード}[turbo mode]では，チップは
1個（または数個）のコアだけを動かし，他を停止して，動作中のコアの電圧と
クロック周波数を通常の値より上げる．$N$ 個のコア分のバジェット全体を使えば
周波数は $\sqrt[3]{N}$ 倍---4コアなら 1.59 倍---まで上げられ，バジェット
の分割で失われた周波数（\cref{eq:l22-freq}）をちょうど取り戻せることになる．
\cref{tab:l22-turbo} は，講義が例として用いる 2013 年の2つのプロセッサを
示す．その設計電力はパッケージ全体のものであり，パッケージには統合
グラフィックスとアンコアも含まれるので，コアだけのバジェットではない．

\begin{table}[htbp]
  \centering\small
  \caption{2つのクアッドコア Intel Core プロセッサの通常時とターボ時の
    クロック周波数．電力の倍率は，ターボ周波数と通常周波数の比の3乗である．}
  \label{tab:l22-turbo}
  \begin{tabular}{@{}lccccc@{}}
    \toprule
    プロセッサ & 設計電力 & 通常 & ターボ & 周波数 & 電力 \\
    \midrule
    i7-4702MQ（ノート） & \SI{37}{W} & \SI{2.2}{GHz} & \SI{3.2}{GHz} & $1.45\times$ & $3.1\times$ \\
    i7-4771（デスクトップ） & \SI{84}{W} & \SI{3.5}{GHz} & \SI{3.9}{GHz} & $1.11\times$ & $1.38\times$ \\
    \bottomrule
  \end{tabular}
\end{table}

どちらのプロセッサも，4コア分のバジェット全体を1つのコアに使ってはいない．
チップを制限するのは総電力だけではなく温度である．ターボモードでは電力が
1つのコアの面積に集中し，その熱の一部しか周囲の停止したコアへ広がらない
ので，動作中のコアは通常の4倍の電力を消費する前に最高温度に達する．ノート
用のプロセッサは通常はコアを低い電力で動かしているので，1つのコアなら電力
をなお約3倍に上げられる．より良い冷却を前提に設計されたデスクトップ用の
プロセッサは，通常の動作でもすでにコアを高い電力で，1つのコアの面積が放熱
できる限界に近いところで動かしており，ターボモードから得られるものは
少ない．

\section{SMT，コア，チップ}
\label{sec:l22-os}

最後の課題はオペレーティングシステムに関わる．\source{Lesson 22，動画
16--17．}メニーコアシステムは第18章のあらゆる種類の並列性を同時に使う．
1枚のボードには複数のチップが載り，各チップには複数のコアがあり，各コアは
SMT で複数のスレッドを走らせる．オペレーティングシステムは，ハードウェア
スレッドのそれぞれを，1つのソフトウェアスレッドを走らせることのできる
\term[ろんりぷろせっさ]{論理プロセッサ}[logical processor]として見る．
例えば，8個のコアとコア当たり2つの SMT スレッドを持つチップが2つあれば，
32 個の論理プロセッサとして見える．

しかし論理プロセッサは等価ではない．同じコアの上の2つのスレッドはその
パイプラインとキャッシュを共有し，それぞれは専用のコアで走る場合よりずっと
遅くなる（\cref{sec:l18-mt-performance}）．同じチップの上のスレッドは最終
レベルキャッシュ，電力バジェット，メモリインタフェースを共有する．異なる
チップの上のスレッドはこれらを何も共有しないが，チップ間のより遅いリンクを
通じて通信する．すべての論理プロセッサを同一とみなすオペレーティング
システムは，3つのスレッドを論理プロセッサ 0，1，2 で開始しかねない．これは
2つのスレッドを1つのコアに置き，もう一方のチップを遊ばせてしまう
（\cref{fig:l22-placement}a）．トポロジを考慮するオペレーティングシステム
は，スレッドをまずチップに，次に各チップのコアに分散させ，あるコアの2つ目
の SMT スレッドを使うのはすべてのコアが埋まってからである
（\cref{fig:l22-placement}b）．

\begin{figure}[htbp]
  \centering
  % Two chips with two cores each and two SMT threads per core give eight
% logical processors (0-7).  Three threads A, B, C placed (a) on logical
% processors 0-2 without regard to the topology, (b) on separate chips and
% cores.
\begin{tikzpicture}[x=1cm,y=1cm, font=\scriptsize\sffamily,
  chip/.style={draw, fill=ln-box, rounded corners=2pt},
  core/.style={draw=ln-muted, fill=white},
  lp/.style={draw, minimum size=0.42cm, inner sep=0pt, font=\footnotesize\sffamily\bfseries},
  busy/.style={fill=ln-accent!25},
  num/.style={font=\tiny\sffamily, text=ln-muted, inner sep=1.5pt},
  paneltitle/.style={font=\footnotesize\sffamily, anchor=north}]
  \foreach \panel/\px/\captext in {0/0/{\iflangja{(a) トポロジを考慮しない配置}{(a) topology-unaware}},
                           1/5.95/{\iflangja{(b) トポロジを考慮した配置}{(b) topology-aware}}} {
    \begin{scope}[xshift=\px cm]
      \foreach \chipi in {0,1} {
        \pgfmathsetmacro{\chx}{\chipi*2.8}
        \draw[chip] (\chx,0) rectangle (\chx+2.55,1.7);
        \node[anchor=north west, inner sep=2pt] at (\chx,1.7) {\iflangja{チップ}{chip}~\chipi};
        \foreach \corei in {0,1} {
          \pgfmathsetmacro{\cox}{\chx+0.15+\corei*1.2}
          \pgfmathtruncatemacro{\coid}{2*\chipi+\corei}
          \draw[core] (\cox,0.1) rectangle (\cox+1.05,1.33);
          \node[anchor=north west, inner sep=1.5pt, text=ln-muted] at (\cox,1.33) {\iflangja{コア}{core}~\coid};
          \foreach \smt in {0,1} {
            \pgfmathtruncatemacro{\lpid}{4*\chipi+2*\corei+\smt}
            \node[lp] (lp-\panel-\lpid) at (\cox+0.3+\smt*0.45,0.7) {};
            \node[num, below] at (lp-\panel-\lpid.south) {\lpid};
          }
        }
      }
      \node[paneltitle] at (2.675,-0.05) {\captext};
    \end{scope}
  }
  % thread placements
  \foreach \panel/\lpid/\thr in {0/0/A, 0/1/B, 0/2/C, 1/0/A, 1/4/B, 1/2/C} {
    \node[lp, busy] at (lp-\panel-\lpid) {\thr};
  }
\end{tikzpicture}

  \caption{2個のコアとコア当たり2つの SMT スレッドを持つ2つのチップの
    上の，3つのスレッド A，B，C．小さな数字は論理プロセッサである．
    (a) スレッド A と B がコア 0 を共有する．(b) どのスレッドも専用の
    コアを持つ．}
  \label{fig:l22-placement}
\end{figure}

\begin{note}
  分散させるのが最良なのは独立なスレッドの場合である．多くのデータを共有
  するスレッドは，データが1つの LLC に留まる同じチップの上でより速く通信
  する．どちらを選ぶにせよ，オペレーティングシステムがトポロジを知っている
  必要がある．
\end{note}

\section{課題と解決策}
\label{sec:l22-summary}

\cref{tab:l22-challenges} はこの章の課題とその解決策をまとめたものである．
\source{Lesson 22，動画 18--19．}いずれの課題も，少数のコアには十分な機構
がコア数とともに成長しないことから生じる．

\begin{table}[htbp]
  \centering\small
  \caption{メニーコアの課題とその解決策．}
  \label{tab:l22-challenges}
  \begin{tabular}{@{}>{\raggedright\arraybackslash}p{0.27\linewidth}>{\raggedright\arraybackslash}p{0.35\linewidth}>{\raggedright\arraybackslash}p{0.31\linewidth}@{}}
    \toprule
    課題 & 原因 & 解決策 \\
    \midrule
    コヒーレンストラフィック & 共有データへの書き込みが増える．バスは一度に
      1つの転送しか運べない & ネットワークオンチップ（メッシュ），
      ディレクトリ方式 \\
    オフチップトラフィック & キャッシュミスが増える．ピン数はコア数より
      ゆっくりとしか増えない & 分散 LLC \\
    ディレクトリの大きさ & メモリブロックごとに1つのエントリ &
      オンチップの部分ディレクトリ \\
    電力バジェット & 1つのバジェットをより多くのコアで分け合う &
      少数のコアだけが動くときのターボモード \\
    スレッドの配置 & SMT スレッド，コア，チップは等価でない &
      トポロジを考慮するオペレーティングシステム \\
    \bottomrule
  \end{tabular}
\end{table}

\begin{keypoint}[章のまとめ]
  メニーコアプロセッサは，第18章から第21章の設計をそのまま大きくしたもの
  にはできない．コヒーレンストラフィックはバスの容量を超えるので，コアは
  ネットワークオンチップ---ふつうはタイルのメッシュ---で接続され，
  ディレクトリ方式を用いる．その容量はバスと違ってコア数とともに増え，
  チップを横切るトラフィックでは $\sqrt{N}$ で，隣り合うタイルの間の
  トラフィックでは $N$ で増える．分散 LLC は各タイルに1つのスライスを
  与え，スライスはキャッシュインデックスで，あるいは局所性のためにページ
  番号で選ばれる．コヒーレンスディレクトリは，各ブロックの LLC スライスの
  タイルに置かれるオンチップの部分ディレクトリになる．共有された電力
  バジェットはどのコアの周波数も $\sqrt[3]{N}$ 分の1にし
  （\cref{eq:l22-freq}），ターボモードは逐次プログラムについてそれを部分的
  に取り戻す．そしてオペレーティングシステムは，どの論理プロセッサがコアや
  チップを共有しているかを知らなければならない．
\end{keypoint}

\end{document}
